\documentclass[]{article}

\usepackage[margin=1in]{geometry}
\usepackage[utf8]{inputenc} % allow utf-8 input
\usepackage[T1]{fontenc}    % use 8-bit T1 fonts
\usepackage{hyperref}       % hyperlinks
\hypersetup{
    colorlinks = true,
    citecolor = blue,
    linkcolor = red
}
\usepackage{url}            % simple URL typesetting
\usepackage{booktabs}       % professional-quality tables
\usepackage{amsfonts}       % blackboard math symbols
\usepackage{nicefrac}       % compact symbols for 1/2, etc.
\usepackage{microtype}      % microtypography
\usepackage{xcolor}         % colors

\usepackage{enumitem}
\setlist[itemize]{leftmargin=0.75em}

\usepackage{color}
\usepackage{float}
\usepackage[style=base]{caption}
\usepackage{subcaption}
\usepackage{tikz}
\usepackage{tikz-cd}
\usepackage{pgfplots}
\usepackage{abstract}
\usepgfplotslibrary{groupplots}
\pgfplotsset{compat=1.18}

\usepackage{amsmath,amssymb,amsthm,mathtools,bbm}
\usepackage{physics}
\usepackage{bbold}
\usepackage{dsfont}
\usepackage{mathrsfs}
\usepackage[ruled,vlined,linesnumbered]{algorithm2e}
\usepackage[capitalize]{cleveref}
\usepackage{xfrac}
\usepackage{graphicx}
\usepackage{authblk}
\usepackage{natbib}
\usepackage{stmaryrd}

\graphicspath{{Figures/}}

\usepackage{times}
\usepackage{comment}

\usepackage{enumitem}
\setlist[enumerate]{leftmargin=1.5em}
\setlist[itemize]{leftmargin=1.5em}
\usepackage{autonum}

\makeatletter
\renewenvironment{proof}[1][]
{%
  \par\noindent{\bfseries\upshape 
    \proofname\if\relax\detokenize{#1}\relax\else\ of #1\fi:\space}%
}
{\qedsymbol}
\makeatother

\input{additional-math.tex}

\newcommand{\defeq}{\mathrel{\mathop:}=}

% % % % END CUSTOM SETTINGS % % % % % % 

%%%%%%%%%%%%%%%%%%%%%%%%%%%%%%%%
% THEOREMS
%%%%%%%%%%%%%%%%%%%%%%%%%%%%%%%%
\theoremstyle{plain}%
\newtheorem{theorem}{Theorem}
\newtheorem{proposition}[theorem]{Proposition}
\newtheorem{lemma}[theorem]{Lemma}
\newtheorem{corollary}[theorem]{Corollary}
\theoremstyle{definition}%
\newtheorem{definition}{Definition}
\newtheorem{assumption}{Assumption}
\theoremstyle{remark}%
\newtheorem{remark}{Remark}

\crefname{assumption}{Assumption}{Assumptions}

\renewcommand\Authfont{\fontsize{11}{14.4}\selectfont}
\renewcommand\Affilfont{\fontsize{10}{10.8}\itshape}

\begin{document}

\title{}

\section{Fixed Temperature Spherical $p$-spin}
We consider the Spherical version of the $p$-spin model, characterized by the following Hamiltonian:
\begin{equation}\label{eq:H_pspin_spherical}
    \mathcal{H}_J(\vect{s}) = -\sum_{i_1<\dots<i_p}J_{i_1\dots i_p}s_{i_1}\dots s_{i_p}\,,\quad \vect{s}\in S_N
\end{equation}
where $J_{i_1\dots i_p}$ is a random Gaussian tensor where each entrance is i.i.d. with mean $\mathbb{E}[J_{i_1\dots i_p}]=\frac{J_0 p!}{N^{p-1}}$ and variance $\mathbb{E}[J^2_{i_1\dots i_p}] - \left(\mathbb{E}[J_{i_1\dots i_p}]\right)^2=\frac{p!}{N^{p-1}}$.
This variance makes the model the same as the one considered in \cite{alaoui_shattering_2024}.

We denote the Gibbs measure $\mu_{\beta}(\vect{s}) = \frac{1}{Z_{\beta}}\exp(-\beta\mathcal{H}_J(\vect{s})) $, with $Z_{\beta}$ the normalization constant.

We use the following definitions for $\beta_d$ and $\beta_K$:
\begin{equation}
    \beta_d(p) = \sqrt{\frac{(p-1)^{p-1}}{p (p-2)^{p-2}}}\,,\quad 
    \beta_K(p)^2 = \inf_{q \in [0,1]} \qty(\frac{1}{q^p} \log\qty(\frac{1}{1-q}) - \frac{1}{q^{p-1}}) \,,
\end{equation}
which for large values of $p$ scale as
\begin{equation}\label{eq:betad-betac-defs}
    \beta_d(p) = \sqrt{e} + O(p^{-1}) \,, \quad \beta_K(p) = \sqrt{\log p} \qty(1 + o_p(1)) \,.
\end{equation}

Our primary objective is to determine $\beta_d^+$, defined as the highest value of $\beta$ for which the system exhibits fast mixing from the worst case initial condition.

\subsection{Bound on $\beta_d^+$ from the monotonicity of the Franz-Parisi potential}

To establish a lower bound on $\beta_d^+$, we follow the proof technique developed in \cite{alaoui_shattering_2024}, applying it to the two temperature Franz Parisi potential. Given an equilibrium configuration from $\vect{s} \sim \mu_{\beta^\prime}$, the two temperature Franz Parisi potential is defined as
\begin{equation}\label{eq:FP-twotemp}
    \mathcal{F}_{\beta,\beta^\prime}(q) :=
    \lim _{\varepsilon \downarrow 0} \operatorname{p-lim}\limits_{N \to \infty} 
    \log 
    \int e^{-\beta H_N\left(\boldsymbol{\vect{s}}^{\prime}\right)} 1_{\left(\vect{s}, \vect{s}^{\prime}\right) \in \mathcal{S}_N(q, \varepsilon)} \mathrm{d} 
    \mu_0\left(\boldsymbol{\vect{s}}^{\prime}\right) \,,
\end{equation}

In order to have contiguity, we consider the case in which we plant at $\beta^\prime \leq \beta_K$. Using \cite[Proposition 3.8]{alaoui_shattering_2024} with 
\begin{equation}
    \vect{G} = \frac{\beta'}{N^{\frac{p-1}{2}}}\vect{s}^{\otimes p} + \vect{W}\,,
\end{equation}
the Franz-Parisi potential in \cref{eq:FP-twotemp} can be written as
\begin{equation}
    \mathcal{F}_{\beta,\beta'}(q) = F_{\beta}(\xi_q) + \beta\beta'q^p + \frac{1}{2}\log(1-q^2)\,,
\end{equation}
with $F_{\beta}(\xi_q) $ the quenched free entropy of the mixed model with covariance $\xi_q$ defined as
\begin{equation}\label{eq:xi_q}
    \xi_q(x) = \left(q^2+\left(1-q^2\right) x\right)^p-q^{2 p} \,.
\end{equation}

By following a similar argument as in \cite[Proposition 3.6]{alaoui_shattering_2024} we can prove the following Lemma.

\begin{lemma}\label{lemm:lowerbound-derivative}
There exists a strictly positive constant $C$ such that for $q\in[1-\frac{1}{p},1)$, the following inequality holds
\begin{equation}\label{eq:Fder}
    \dv{\mathcal{F}_{\beta,\beta'}}{q} \geq -\frac{(C\beta + 1)q}{1-q^2} + \beta\beta'pq^{p-1} \,.
\end{equation}
\end{lemma}

\begin{remark}
Compared to \cite{alaoui_shattering_2024}, we had to improve the range of validity to $q > 1-\frac{1}{p}$, and we show how to get this improvement in the next subsection.
\end{remark}


% \begin{lemma}\label{lem:beta-d-plus}
% For all $\varepsilon > 0$ and all $\beta^\prime \leq \beta_K$ there exists $p_0 \in \mathbb{N}$ such that for all $p \geq p_0$ we have
% \begin{equation}
%     \abs{\beta_d^+ - \frac{e}{2 \beta^\prime} } < \varepsilon \,.
% \end{equation}
% \end{lemma}

\begin{lemma}\label{lem:beta-d-plus}
Take $\beta'\leq\beta_K$ and $\beta = o_p(1)$. Then, for all $\delta > 0$ 
there exists $p_0(\delta) \in \mathbb{N}$ and $\alpha^\star(\delta) < 1$ such that for all $p \geq p_0(\delta)$ and all $q \in [1-\frac{1}{p}, 1- \frac{\alpha^\star(\delta)}{p}]$, if $\beta\beta' > \frac{e}{2} + \delta$ we have that 
\begin{equation}
    \frac{ \dd }{ \dd q} \mathcal{F}_{\beta,\beta'} (q) > 0 \,.
\end{equation}
\end{lemma}

\begin{remark}
If we take $\beta' = \beta_K$, the smallest value of $\beta$ for which the lemma holds is $\frac{e}2{(\log p)^{-\frac{1}{2}}}$, which coincides with the value of $\beta_d^+$ one gets from using a Replica Symmetric ansatz for the Franz-Parisi potential. We detail such derivation in the next section.
\end{remark}

\begin{proof}[\Cref{lemm:lowerbound-derivative}]
We begin by analyzing the series coefficients $\xi_{q,k}$ of $\xi_q(x)$ as defined in \cref{eq:xi_q}. These coefficients are given by
\begin{equation}\label{eq:xi_q_k}
    \xi_{q,k} = \begin{cases}
        \binom{p}{k} \qty(1 - q^2)^k q^{2 (p-k)} & \text{ for } k \in \llbracket 1,p \rrbracket \\ 
        0 & \text{ for } k = 0
    \end{cases} \,.
\end{equation}

% In \cite{alaoui_shattering_2024}, it is used the fact that for $q \in [1 - \frac{1}{2p},1)$ one can show that $\xi_{q,k}\leq 1/k!$ and this allows to show that the derivative of the free energy is upper bounded by a universal constant. This argument is not directly generalizable if one takes $q<1 - \frac{1}{2p}$, and one has to compute a different bound. 

Let's start by considering the $q$ for which $\xi_{q,k}$ is maximized. Specifically
\begin{equation}
    q^{\mathrm{max}}_{k} := \argmax\limits_{q \in [1-\frac{1}{p}, 1)} \xi_{q,k} =
    \begin{cases}
        \sqrt{1-\frac{1}{p}} & \text{ for } k = 1 \\
        1 - \frac{1}{p} & \text{ for } k > 1
    \end{cases}
\end{equation}
for $p \in \mathbb{N}$ and $p \geq 2$.

By substituting $q^{\mathrm{max}}_{k}$ into the formula \cref{eq:xi_q_k}, for $k>1$ we get 
\begin{equation}
    \xi^*(j) := p^{-2 p} (p-1)^{2 p-2 j} (2 p-1)^j \binom{p}{j} < p^{-2 p} p^{2 p-2 j} (2p)^j \binom{p}{j}\,,
\end{equation}
% and using the fact that $\binom{p}{j} \leq p^j/j!$ we get
% \begin{equation}
%     \xi^*(j) < \frac{2^j}{j!}\,.
% \end{equation}
We can further bound this expression as follows:
\begin{align}
    \xi^*(k) &< p^{-2p}p^{2p-2k}(2p)^k \frac{p^k}{k!} = p^{-2p}p^{2p-k}2^k \frac{p^k}{k!} = \frac{2^k}{k!}.
\end{align}
In the inequality, we used the bound $\binom{p}{k} \leq \frac{p^k}{k!}$.

Using the proof technique of \cite[Proposition 3.6]{alaoui_shattering_2024} we need to check that 
\begin{equation}
    \sum_{j=1}^p \sqrt{\xi_{q,j}} \sqrt{ \log(j)}
\end{equation}
is finite for any value of $p$. This is the case for example if we can show that $\xi^*(j) < \frac{1}{j^2}$ for all $j$ large enough.

Thus, by calling $A(j) = \frac{2^j j^2}{j!}$ we have
\begin{equation}
    \xi^*(j) < \frac{A(j)}{j^2}\,,
\end{equation}
and since $A(j) < 1\;,\,\forall j>7$ we get 
\begin{equation}
    \xi^*(j) < \frac{1}{j^2}\;,\,\forall j>7\,.
\end{equation}

This bound, combined with the analysis from \cite[Proposition 3.6]{alaoui_shattering_2024}, establishes the desired lower bound on the derivative of the Franz-Parisi potential as stated in equation \cref{eq:Fder}.
\end{proof}

\begin{proof}[\Cref{lem:beta-d-plus}]
%We look at the first derivative of $\mathcal{F}_{\beta,\beta^\prime}(q)$, and we want to characterize the minimum value of $\beta$ for which there exists a range in $q$ the derivative is positive. 
%Indeed, this condition is enough to establish the non-monotonicity of the Franz-Parisi potential.
We start from the results of \Cref{lemm:lowerbound-derivative} and we parametrize $q = 1 - \frac{\alpha}{p}$ with $\alpha \in (0,1]$.% and $\beta = \tilde{\beta}(\log p)^{-\frac{1}{2}}$ with $\tilde{\beta}$ a positive constant.

Then the RHS of \cref{eq:Fder} then becomes
\begin{equation}
    \beta\beta' p \left(1-\frac{\alpha }{p}\right)^{p-1} -
    \frac{p \left(1-\frac{\alpha }{p}\right) \left(\beta C+1\right)}{2 \alpha} \qty( 1 + \frac{\alpha}{2 p} + \frac{\alpha^2}{4 p^2} + o_p(p^{-2}) ) \,.
\end{equation}

Keeping only the first order, if we assume $\beta = o_p(1)$ this becomes
\begin{equation}
    -\frac{p}{2 \alpha} + e^{-\alpha}\beta\beta' p + o_p(p)
\end{equation}

If we now choose $\beta$ such that $\beta\beta' = \frac{e}{2} + \delta$, we can study the function $h_{\delta} : \alpha \mapsto -\frac{1}{2 \alpha} + \frac{e^{1-\alpha}}{2} + \delta e^{-\alpha}$.

Putting this equation to zero gives only one solution in $\alpha \in (0,1]$. We call this solution $\alpha^\star$ and we have that 
\begin{equation}
    \alpha^\star(\delta) = -W_0\qty(-\frac{1}{e+ 2 \delta})
\end{equation}
where $W_0$ is the principal solution Lambert function.

Consequently, for any $\alpha \in [\alpha^*(\delta), 1]$, there exists a sufficiently large $p_0$ such that for all $p > p_0$, we have:
\begin{equation}
    \frac{d\mathcal{F}_{\beta,\beta'}}{dq}(q) > 0,
\end{equation}
which completes the proof.
\end{proof}

% We require the right-hand-side of \cref{eq:Fder} to be positive. For $\beta$ this means that
% \begin{equation}\label{eq:beta-lower}
%     \beta \geq \frac{1}{- C + (1-q^2)p\beta^\prime q^{p-2}} \geq \frac{1}{(1-q^2)p \beta^\prime q^{p-2}} \,,
% \end{equation}
% where we used the fact that $C$ is strictly positive.

% For any fixed $\beta^\prime$ the previous function is strictly positive for $q\in[1-\frac{1}{p},1)$ for any value of $p$. 
% % Then the value of $q^*$ which maximizes the denominator corresponds to the minimum value $\beta(q^*)$ for which the non-monotonicity condition is satisfied.
% We parametrize $q(\alpha)=1-\frac{\alpha}{p}$, where then $\alpha\in(0,1]$ to get
% \begin{equation}
%     \beta \geq \frac{1}{\alpha(2-\frac{\alpha}{p})\beta^\prime(1-\frac{\alpha}{p})^{p-2}} \geq \frac{1}{2\alpha \beta^\prime (1-\frac{\alpha}{p})^{p-2}} \,,
% \end{equation}
% where the last inequality follows because $\alpha \in (0,1]$.
% The sequence $a_n = \alpha^{-1}\qty(1-\frac{\alpha}{n})^{-n+2}$ is a monotonically increasing bounded above by $e^{\alpha}/\alpha$ and thus converges to this value. 
% By the diefinition of the limilt $\forall \varepsilon > 0$ $\exists p_0$ such that for $p \geq p_0$
% \begin{equation}
%     2 \beta \beta^\prime \geq a_p \geq \frac{e^{\alpha}}{\alpha} - \varepsilon \geq e -\varepsilon \,,
% \end{equation}
% where it follows because $\sfrac{e^{\alpha}}{\alpha} \geq e$.



% \begin{equation}
%     2 \sqrt{\log p} \beta(\alpha) > \frac{e^\alpha}{\alpha} - \epsilon 
% \end{equation}
% Then we can choose $\epsilon$ such that for every $\alpha$
% \begin{equation} 
%     \beta(\alpha) \geq \frac{e}{2 \sqrt{\log p}}
% \end{equation}
% thus proving that in the large p limit we have the same expression as the RS prediction found in the following section.

\newpage
\subsection{Computation of $T_d^+$ from following states (RS)}
The RS Franz-Parisi potential can be shown to be
\begin{equation}\label{eq:FP_RS_spherical}
    -2\beta F_\beta(m,q) = 1 + \log 2\pi + \beta^2(1-q^p) + \log(1-q) + \frac{q-m^2}{1-q} + 2\beta J_0 m^p\,.
\end{equation}
where $J_0 = \beta' $, with $\beta'$ the inverse temperature of the planted solution, $m$ is the overlap with the planted solution and $q$ is the replica overlap.

\subsubsection{Finite $p$}
We first use the state following formalism from \citep{krzakala_2010,sun_2012} to get an estimate of the $T_d^+$ for the spherical case at a finite value of $p$.
We consider the planted version of the model, with the planting temperature fixed at the Kauzmann temperature $T_K$, so that we have contiguity to the random version of the model.
We start by running the fixed point equations at temperature $T = T_K$ to find the local minima with $m > 0$.
Then, we increase the temperature and follow the minimum until it vanishes, and we call $T_d^+$ the temperature value at which this happens.
For finite values of p (for $p > 10^4$, the numerical routines start to become unstable), we can look at the scaling of $T_d^+$ with p and it turns out to be consistent with the one found in the previous section, as we will show in the plots in the next section.
%In Figure \ref{fig:spherical_Tdplus_Td_Tk} we plot these values, along with the values of $T_d$ and $T_K$.

\subsubsection{Large $p$ limit}
We can rewrite the fixed-point equations of \eqref{eq:FP_RS_spherical} as following
\begin{equation}\label{eq:fixedpoints}
\begin{cases}
    1-q = \frac{m^{2-p}}{\beta p J_0}\,, \\
    (1-q)^2 = \frac{(q-m^2)q^{1-p}}{p\beta^2}\,.
\end{cases}
\end{equation}
In the large $p$ limit the Kauzmann transition is at $\beta_K \sim \sqrt{\log p}$ \footnote{We use the symbol $\sim$ to indicate that two expressions are of the same order in $p$.}\cite{kurchan1993barriers} and, following the numerical results, we consider $\beta = \frac{C_+}{\sqrt{\log p}}$, with $C_+$ is a constant to be determined.

Then from the first fixed point we get $1-q \sim (pm^p)^{-1}$. Numerically, one can verify that $m^p$ converges for $p$ large to a positive constant, such that we can say $q = 1 - \frac{\alpha_q}{p} + o(\frac{1}{p})$, and thus $q^p\to e^{-\alpha_q}$.

Putting this scaling in the second fixed point, we get
\begin{equation}
    \frac{\alpha_q}{p} \sim \sqrt{\frac{\log p}{p}(q-m^2)}\,,
\end{equation}
which in turn implies $q-m^2 = \frac{\gamma}{p\log p}$. By a simple argument, one can see that this means $m = 1 - \frac{\alpha_m}{p}= 1 - \frac{\alpha_q}{2p} + o(\frac{1}{p})$ and thus that $m^p\to e^{-\alpha_q/2}$.

Directly by imposing these scaling back into \eqref{eq:fixedpoints} one finds the following relationships between the constants:
\begin{equation}\label{eq:Cplusalpha}
    \begin{cases}
        C_+\alpha_q = e^{\frac{\alpha_q}{2}}\,, \\
        C_+^2\alpha_q^2 = \gamma e^{\alpha_q}\,.
    \end{cases}
\end{equation}
Using the first equation to specify $C_+$ in terms of $\alpha_q$ and substituting in the second equation we find $\gamma = 1$ and we are left with one equation in two variables.

We derive the last equation by looking at the Hessian of the free energy $F_{\beta}(m,q)$. Indeed we expect that $T_d^+$ is the temperature at which the local minima disappears, and thus it is also where the smaller eigenvalue of the Hessian becomes zero. 
The eigenvalues can be computed and the smaller one is given by

\begin{equation}
    \begin{cases}
        \lambda_{\rm min} = -\frac{1}{4 \beta  m^2 (q-1)^3 q^2}\left(A - \sqrt{A^2 + B}\right)\,, \\
        \begin{aligned}
        A = 2 \beta  J_0 (p-1) p (q-1)^3 q^2 m^p+2 m^4 q^2 -m^2 \left(\beta ^2 (p-1) p (q-1)^3 q^p- q^2 (q (2 q-5)+1)\right) \,, 
        \end{aligned}\\
        \begin{aligned}
        B = 8&m^2 (q-1)^2 q^2 \Bigg(\beta  J_0 (p-1) p (q-1) m^p \left( q^2 \left(-2 m^2+q+1\right)+\beta ^2 (p-1) p (q-1)^3
   q^p\right)\\ &+m^2 \left(\beta ^2 (p-1) p (q-1)^3 q^p+ (q+1) q^2\right)\Bigg)\,.
        \end{aligned}
    \end{cases}
\end{equation}
One can check that $B = o(A)$ and thus it is easy to convince ourselves that $B = 0 \Rightarrow\lambda_{\min} = 0$. At the first order in $B$, this is equivalent to asking 
\begin{equation}
    m^2q^2(1+q) + m^pp^2(q-1)q^2\beta J_0(1-2m^2 + q) = 0 \,.
\end{equation}
After some trivial computations, one arrives to the following formula:
\begin{equation}
    C_+\alpha_q^2 = 2e^{\frac{\alpha_q}{2}} \,,
\end{equation}
which combined with the one coming from the fixed point equations leads to $\alpha_q = 2$, and thus $\alpha_m=1$ and $C_+ = \frac{e}{2}$.

We verify these scalings numerically; in Fig. \ref{fig:spherical2_Tdplus_Td_Tk_log} for $T_d^+$ and in Fig. \ref{fig:spherical2_m_q} for $m$ and $q$.
\begin{figure}
    \centering
    \includegraphics[width=0.75\linewidth]{Figures/spherical2_Tdplus_TkTd_log.png}
    \caption{Numerical estimation of $T_d^+$ from state following, superimposed with the asymptotic scalings.}
    \label{fig:spherical2_Tdplus_Td_Tk_log}
\end{figure}
\begin{figure}
    \centering
    \includegraphics[width=1\linewidth]{Figures/spherical2_m_q.png}
    \caption{Check of the asymptotic scalings of $m$ and $q$.}
    \label{fig:spherical2_m_q}
\end{figure}

Note: Another way to find this scalings is to look at
\eqref{eq:Cplusalpha} and see what is the smallest $C_+$ for which there is a solution. 

 \bibliographystyle{alpha} % We choose the "plain" reference style
 \bibliography{refs}

\end{document}